\documentclass[12pt]{article}

\usepackage{amsmath, amssymb, amsthm}
\usepackage{mathrsfs, mathtools}
\usepackage{geometry}
\usepackage{hyperref}
\usepackage{natbib}
\usepackage{enumitem}

\geometry{margin=1.25in}

\newtheorem{theorem}{Theorem}[section]
\newtheorem{lemma}[theorem]{Lemma}
\newtheorem{corollary}[theorem]{Corollary}
\newtheorem{proposition}[theorem]{Proposition}
\newtheorem{definition}[theorem]{Definition}
\newtheorem{remark}[theorem]{Remark}

\DeclareMathOperator{\esssup}{ess\,sup}
\newcommand{\norm}[1]{\left\|#1\right\|}
\newcommand{\abs}[1]{\left|#1\right|}
\newcommand{\E}{\mathbb{E}}
\newcommand{\Ec}[2]{\mathbb{E}\!\left[\,#1\,\middle|\,#2\,\right]}
\newcommand{\F}{\mathscr{F}}
\newcommand{\TV}[1]{\left\|#1\right\|_{\mathrm{TV}}}
\newcommand{\Prob}{\mathbb{P}}

\title{The Blackwell--Dubins Theorem on the Merging of Opinions:\\
       Discrete and Continuous Time}
\author{Analytical Exposition}
\date{}

\begin{document}
\maketitle
\tableofcontents
\bigskip

%----------------------------------------------------------------------
\section{Introduction}
%----------------------------------------------------------------------

A foundational question in Bayesian epistemology and statistical decision
theory is: if two agents hold different prior beliefs about a stochastic
process but observe the same data indefinitely, will their probability
assessments eventually converge?  \citet{blackwell1962} give a precise and
general affirmative answer under the single hypothesis that the two priors
are mutually absolutely continuous.  Their result is now called the
\textbf{merging-of-opinions theorem}, and it occupies a central place in
Bayesian learning, rational expectations theory, and the ergodic theory of
Markov chains.

The theorem's content can be stated informally as follows.  Let $P$ and $Q$
be two probability measures on the space of infinite sequences (or infinite
paths) of observations.  If $P$ is absolutely continuous with respect to $Q$,
then---as the observed history grows without bound---the \emph{conditional}
distributions of the future path assigned by $P$ and $Q$ merge in total
variation with probability one under both measures.  No matter how different
the two agents' initial beliefs were, they will, almost surely, make
identical probability assessments about all future events once enough data
have been observed.

Section~\ref{sec:setup} develops the measurable-space framework and recalls
the total variation metric.  Section~\ref{sec:RN} introduces the
Radon--Nikodym likelihood-ratio process and its martingale properties.
Section~\ref{sec:BD} states and proves the Blackwell--Dubins theorem in
discrete time.  Section~\ref{sec:CT} extends the analysis to continuous
time on $[0, +\infty)$, where the Girsanov--Cameron--Martin theorem plays
the role of the Radon--Nikodym derivative.
Section~\ref{sec:kakutani} discusses Kakutani's dichotomy, which characterises
exactly when mutual absolute continuity holds for product measures, and
hence when merging is guaranteed.  Section~\ref{sec:applications} discusses
applications to Bayesian learning, rational expectations, and the ergodic
theory of diffusions.  Section~\ref{sec:related} collects related results
and extensions.

%----------------------------------------------------------------------
\section{Setup and Notation}
\label{sec:setup}
%----------------------------------------------------------------------

\subsection{The canonical sequence space}

Let $(S, \mathscr{S})$ be a measurable space (the \emph{signal space}).  Set
$\Omega = S^{\mathbb{N}}$, the space of all sequences $\omega =
(x_1, x_2, \ldots)$ with $x_n \in S$.  Equip $\Omega$ with the product
$\sigma$-algebra $\mathscr{F} = \mathscr{S}^{\otimes\mathbb{N}}$.

For each $n \geq 1$ define the finite-horizon $\sigma$-algebra
\[
\mathscr{F}_n = \sigma(x_1, \ldots, x_n),
\]
so that $\mathscr{F}_1 \subseteq \mathscr{F}_2 \subseteq \cdots \subseteq
\mathscr{F}$.  Let $\mathscr{F}_\infty = \sigma\bigl(\bigcup_{n \geq 1}
\mathscr{F}_n\bigr)$.  When $S$ is a Polish space and $\mathscr{S}$ is the
Borel $\sigma$-algebra, $\mathscr{F}_\infty = \mathscr{F}$.

Let $P$ and $Q$ denote two probability measures on $(\Omega, \mathscr{F})$.
Write $P_n = P|_{\mathscr{F}_n}$ and $Q_n = Q|_{\mathscr{F}_n}$ for their
restrictions to the history up to time $n$.

\subsection{Absolute continuity}

\begin{definition}[Absolute continuity on a filtered space]
$P$ is \textbf{absolutely continuous} with respect to $Q$, written $P \ll Q$,
if for every $A \in \mathscr{F}$, $Q(A) = 0$ implies $P(A) = 0$.  They are
\textbf{mutually absolutely continuous} (or \emph{equivalent}), written
$P \sim Q$, if both $P \ll Q$ and $Q \ll P$.

$P$ is \textbf{locally absolutely continuous} with respect to $Q$ if
$P_n \ll Q_n$ for every $n \geq 1$; global absolute continuity
$P \ll Q$ implies local absolute continuity but is strictly stronger.
\end{definition}

\subsection{Total variation distance}

For two probability measures $\mu$ and $\nu$ on a measurable space
$(E, \mathscr{E})$, the \textbf{total variation distance} is
\[
\TV{\mu - \nu}
= \sup_{A \in \mathscr{E}}\abs{\mu(A) - \nu(A)}
= \tfrac{1}{2}\int_E \abs{\frac{d\mu}{d\lambda} - \frac{d\nu}{d\lambda}}\, d\lambda,
\]
where $\lambda$ is any common dominating measure.  Equivalently,
\begin{equation}
\label{eq:TV-formula}
\TV{\mu - \nu} = 1 - \int_E \min\!\left(\frac{d\mu}{d\lambda},
  \frac{d\nu}{d\lambda}\right) d\lambda,
\end{equation}
which takes values in $[0, 1]$.  When $\mu \ll \nu$ with $f = d\mu/d\nu$,
\begin{equation}
\label{eq:TV-RN}
\TV{\mu - \nu} = \E_\nu[(f - 1)^+] = \E_\nu[(1 - f)^+]
= 1 - \E_\nu[\min(f, 1)].
\end{equation}

\subsection{Conditional distributions and the merging question}

The Blackwell--Dubins theorem concerns the conditional distributions of the
\emph{future} given the \emph{past}.  For each $n$ and each $\omega$, let
$P(\cdot\,|\,\mathscr{F}_n)(\omega)$ denote the conditional probability
measure on $(\Omega, \mathscr{F})$ induced by $P$ given the event
$\{x_1 = \omega_1, \ldots, x_n = \omega_n\}$.  Define analogously
$Q(\cdot\,|\,\mathscr{F}_n)(\omega)$.

\begin{definition}[Merging in total variation]
$P$ and $Q$ \textbf{merge in total variation} if
\[
\TV{P(\cdot\,|\,\mathscr{F}_n) - Q(\cdot\,|\,\mathscr{F}_n)} \to 0
\quad \text{almost surely.}
\]
The qualifier ``almost surely'' is with respect to the measure under which
the convergence is asserted, typically $Q$ (and hence also $P$ when $P \ll Q$).
\end{definition}

%----------------------------------------------------------------------
\section{The Likelihood-Ratio Martingale}
\label{sec:RN}
%----------------------------------------------------------------------

\subsection{The Radon--Nikodym process}

Since $P \ll Q$ implies $P_n \ll Q_n$ for every $n$, the
\textbf{Radon--Nikodym derivative} (likelihood ratio) at stage $n$ is
\[
Z_n = \frac{dP_n}{dQ_n}, \qquad Z_n \geq 0 \;\; Q\text{-a.s.},
\qquad \E_Q[Z_n] = 1.
\]
The global condition $P \ll Q$ guarantees that $Z = dP/dQ$ exists on
$(\Omega, \mathscr{F})$, and
\begin{equation}
\label{eq:Zn-tower}
Z_n = \Ec{Z}{\mathscr{F}_n} \quad Q\text{-a.s.}
\end{equation}
That is, $\{Z_n, \mathscr{F}_n\}_{n \geq 1}$ is a \textbf{non-negative,
uniformly integrable $Q$-martingale}.

\begin{lemma}[Martingale convergence]
\label{lem:mart-conv}
The likelihood ratio process $\{Z_n\}$ satisfies:
\begin{enumerate}[label=(\roman*)]
\item $Z_n \to Z_\infty$ $Q$-a.s.\ as $n \to \infty$.
\item $Z_\infty = \Ec{Z}{\mathscr{F}_\infty}$ $Q$-a.s.
\item $Z_n \to Z_\infty$ in $L^1(Q)$, i.e., $\E_Q[|Z_n - Z_\infty|] \to 0$.
\item When $\mathscr{F}_\infty = \mathscr{F}$, $Z_\infty = Z$ $Q$-a.s.,
      and the $L^1$ convergence is a consequence of uniform integrability.
\end{enumerate}
\end{lemma}

\begin{proof}
Parts (i) and (ii) follow from the almost-sure martingale convergence theorem
\citep{doob1953}, since $\{Z_n\}$ is a non-negative martingale and hence
bounded in $L^1(Q)$.  Uniform integrability of $\{Z_n\}$ follows from $Z
\in L^1(Q)$ (via \eqref{eq:Zn-tower} and Jensen's inequality), which gives
(iii).  Part (iv) is the identification of the a.s.\ limit with the terminal
Radon--Nikodym derivative.
\end{proof}

\subsection{The Bayes formula for conditional posteriors}

The following identity relates the conditional measures to $Z_n$:

\begin{lemma}[Conditional Bayes formula]
\label{lem:bayes}
On the set $\{Z_n > 0\}$, for any $A \in \mathscr{F}$,
\begin{equation}
\label{eq:bayes-cond}
P(A\,|\,\mathscr{F}_n) = \frac{\Ec{Z_\infty \cdot \mathbf{1}_A}{\mathscr{F}_n}}{Z_n}
\quad Q\text{-a.s.}
\end{equation}
Equivalently, the Radon--Nikodym derivative of $P(\cdot\,|\,\mathscr{F}_n)$
with respect to $Q(\cdot\,|\,\mathscr{F}_n)$ on $\mathscr{F}$ is
\begin{equation}
\label{eq:cond-RN}
\frac{dP(\cdot\,|\,\mathscr{F}_n)}{dQ(\cdot\,|\,\mathscr{F}_n)}
= \frac{Z_\infty}{Z_n} \quad Q\text{-a.s.\ on } \{Z_n > 0\}.
\end{equation}
\end{lemma}

\begin{proof}
Equation~\eqref{eq:bayes-cond} is an application of the abstract Bayes
theorem: for $B \in \mathscr{F}_n$ and $A \in \mathscr{F}$,
\[
P(A \cap B) = \int_B \Ec{Z_\infty \mathbf{1}_A}{\mathscr{F}_n} dQ
= \int_B \frac{\Ec{Z_\infty \mathbf{1}_A}{\mathscr{F}_n}}{Z_n} Z_n \, dQ
= \int_B \frac{\Ec{Z_\infty \mathbf{1}_A}{\mathscr{F}_n}}{Z_n} dP_n,
\]
which identifies \eqref{eq:bayes-cond}.  Equation~\eqref{eq:cond-RN} follows
by applying \eqref{eq:TV-RN} to the ratio $Z_\infty/Z_n$.
\end{proof}

%----------------------------------------------------------------------
\section{The Blackwell--Dubins Theorem: Discrete Time}
\label{sec:BD}
%----------------------------------------------------------------------

\subsection{Statement}

\begin{theorem}[Blackwell--Dubins, 1962]
\label{thm:BD}
Let $P$ and $Q$ be probability measures on $(\Omega, \mathscr{F})$ with
$P \ll Q$.  Let
\[
d_n = \TV{P(\cdot\,|\,\mathscr{F}_n) - Q(\cdot\,|\,\mathscr{F}_n)}.
\]
Then $d_n \to 0$ almost surely under $Q$ (and hence also under $P$).
\end{theorem}

\subsection{Proof}

The proof combines three ingredients: (a) an expression for $d_n$ in terms
of the likelihood-ratio process; (b) the supermartingale property of
$\{d_n\}$; and (c) $L^1$ convergence of the likelihood ratio.

\paragraph{Step 1: Expression for $d_n$.}

By Lemma~\ref{lem:bayes}, on $\{Z_n > 0\}$ the conditional measure
$P(\cdot\,|\,\mathscr{F}_n)$ is absolutely continuous with respect to
$Q(\cdot\,|\,\mathscr{F}_n)$ with Radon--Nikodym derivative $Z_\infty/Z_n$.
Applying \eqref{eq:TV-RN} gives
\begin{equation}
\label{eq:dn-formula}
d_n
= 1 - \Ec{\min\!\left(\frac{Z_\infty}{Z_n}, 1\right) \cdot Z_n}{\mathscr{F}_n}_{Q}
= \Ec{\left(1 - \frac{Z_\infty}{Z_n}\right)^+}{\mathscr{F}_n}_{Q(\cdot|\mathscr{F}_n)}.
\end{equation}
Equivalently, after multiplying through by $Z_n$ and taking the $Q$-expectation:
\begin{equation}
\label{eq:dn-L1}
2\, d_n\, Z_n
= \Ec{|Z_\infty - Z_n|}{\mathscr{F}_n}_{Q}
\quad Q\text{-a.s.}
\end{equation}

\paragraph{Step 2: $\{d_n\}$ is a $Q$-supermartingale.}

Conditioning on more information cannot increase the difficulty of
distinguishing two hypotheses; it can only reduce it.  Formally,
for any coupling of $P(\cdot|\mathscr{F}_n)$ and $Q(\cdot|\mathscr{F}_n)$,
\[
d_n = \TV{P(\cdot|\mathscr{F}_n) - Q(\cdot|\mathscr{F}_n)}
\geq \TV{P(\cdot|\mathscr{F}_{n+1}) - Q(\cdot|\mathscr{F}_{n+1})}
\quad Q\text{-a.s.}
\]
More precisely, since $P(\cdot|\mathscr{F}_n) = \Ec{P(\cdot|\mathscr{F}_{n+1})}{\mathscr{F}_n}$
and the TV norm is convex,
\begin{equation}
\label{eq:supermg}
\Ec{d_{n+1}}{\mathscr{F}_n} \leq d_n \quad Q\text{-a.s.}
\end{equation}
Thus $\{d_n, \mathscr{F}_n\}$ is a non-negative $Q$-supermartingale bounded
in $[0,1]$.  By the supermartingale convergence theorem
\citep{doob1953}, $d_n \to d_\infty$ $Q$-almost surely for some
$[0,1]$-valued random variable $d_\infty$.

\paragraph{Step 3: The limit is zero.}

From \eqref{eq:dn-L1} and the tower property,
\[
\E_Q[d_n] \leq \tfrac{1}{2}\E_Q[|Z_\infty - Z_n|] \to 0
\]
by the $L^1$ convergence in Lemma~\ref{lem:mart-conv}(iii).  Hence
$d_n \to 0$ in $L^1(Q)$.  Since $d_n \to d_\infty$ $Q$-a.s.\ and
$d_n \to 0$ in probability under $Q$, the two limits must coincide:
$d_\infty = 0$ $Q$-a.s.

Since $P \ll Q$, every $Q$-null set is also $P$-null, so
$d_n \to 0$ $P$-a.s.\ as well. \hfill$\square$

\subsection{Tightness and the role of mutual absolute continuity}

\begin{remark}[One-sided vs.\ mutual absolute continuity]
The theorem requires only $P \ll Q$; it does not require $Q \ll P$.
Under the one-sided condition:
\begin{itemize}
\item Merging holds $Q$-a.s.\ (and hence $P$-a.s.\ since $P \ll Q$).
\item If additionally $Q \ll P$ (i.e., $P \sim Q$), merging holds a.s.\
      under both measures, and both agents learn to agree.
\end{itemize}
\end{remark}

\begin{remark}[Tightness]
The theorem is tight in the following sense.  If $P \perp Q$ (mutual
singularity), then there exists a set $A$ with $P(A) = 1$ and $Q(A) = 0$.
By Lévy's zero-one law, $Q(A|\mathscr{F}_n) \to 0$ and
$P(A|\mathscr{F}_n) \to 1$ almost surely, so $d_n \to 1$ rather
than zero.  The dichotomy between $P \ll Q$ and $P \perp Q$ is thus
qualitatively sharp.
\end{remark}

%----------------------------------------------------------------------
\section{Extension to Continuous Time on $[0, +\infty)$}
\label{sec:CT}
%----------------------------------------------------------------------

\subsection{Filtered probability spaces and local martingales}

Let $(\Omega, \mathscr{F}, \{\mathscr{F}_t\}_{t \geq 0})$ be a filtered
measurable space satisfying the usual conditions (right-continuity of
$\{\mathscr{F}_t\}$ and completeness of $\mathscr{F}_0$).  Let $P$ and $Q$
be two probability measures on $(\Omega, \mathscr{F}_\infty)$ where
$\mathscr{F}_\infty = \sigma(\bigcup_{t \geq 0} \mathscr{F}_t)$.

Denote by $P_t = P|_{\mathscr{F}_t}$ and $Q_t = Q|_{\mathscr{F}_t}$ the
restrictions to the filtration at time $t$.

\begin{definition}[Local absolute continuity]
$P$ is \textbf{locally absolutely continuous} with respect to $Q$ if
$P_t \ll Q_t$ for all $t < \infty$.  When additionally $P \ll Q$ on
$\mathscr{F}_\infty$, one says $P$ is \textbf{globally absolutely continuous}.
\end{definition}

\noindent Note that local absolute continuity does \emph{not} imply global
absolute continuity.  The gap between them is a key subtlety of the
continuous-time theory.

\subsection{The likelihood-ratio process via Girsanov's theorem}

On the canonical Wiener space, $(\Omega, \mathscr{F}, Q)$ is the Wiener
measure so that the coordinate process $W = \{W_t\}$ is a standard
Brownian motion.  Suppose that under $P$ the process $W$ has an additional
drift: $P$-almost surely,
\[
W_t = \widetilde{W}_t + \int_0^t \theta_s\, ds,
\]
where $\widetilde{W}$ is a $P$-Brownian motion and $\{\theta_s\}$ is an
$\{\mathscr{F}_s\}$-adapted process.  The \textbf{Girsanov--Cameron--Martin
theorem} \citep{girsanov1960} states that, if the stochastic exponential
\begin{equation}
\label{eq:stoch-exp}
Z_t = \mathcal{E}\!\left(\int_0^\cdot \theta_s\, dW_s\right)_t
= \exp\!\left(\int_0^t \theta_s\, dW_s - \frac{1}{2}\int_0^t \theta_s^2\, ds\right)
\end{equation}
is a $Q$-martingale (not merely a local martingale), then $P_t \ll Q_t$ for
all $t$ with $dP_t/dQ_t = Z_t$.

The process $Z_t$ is always a non-negative $Q$-local martingale; it is a
true martingale if and only if $\E_Q[Z_t] = 1$ for all $t$.  Sufficient
conditions for this are:

\begin{itemize}
\item \textbf{Novikov's condition} \citep{novikov1972}: $\E_Q\!\left[\exp\!\left(\tfrac{1}{2}\int_0^T \theta_s^2\, ds\right)\right] < \infty$ for all $T$.
\item \textbf{Kazamaki's condition}: the process $\exp\!\left(\tfrac{1}{2}\int_0^\cdot \theta_s\, dW_s\right)$ is a uniformly integrable $Q$-submartingale.
\end{itemize}

\subsection{Global vs.\ local absolute continuity on $[0,+\infty)$}

A central issue on the infinite time horizon is that local absolute
continuity ($Z_t$ a true martingale for all $t$) does \emph{not} ensure
$P \ll Q$ on $\mathscr{F}_\infty$.

\begin{proposition}[Dichotomy at infinity]
\label{prop:dichotomy}
Suppose $Z_t$ defined by \eqref{eq:stoch-exp} is a true $Q$-martingale on
$[0, T]$ for every finite $T$.  Then:
\begin{enumerate}[label=(\roman*)]
\item $Z_t \to Z_\infty$ $Q$-a.s.\ as $t \to \infty$.
\item If $\{Z_t\}$ is uniformly integrable over $[0,\infty)$, then
      $P \ll Q$ on $\mathscr{F}_\infty$ with $dP/dQ = Z_\infty > 0$
      $P$-a.s.
\item If $\{Z_t\}$ is \emph{not} uniformly integrable, then $Z_\infty = 0$
      $Q$-a.s., and $P \perp Q$ on $\mathscr{F}_\infty$.
\end{enumerate}
In case (iii), $P$ and $Q$ become mutually singular at the terminal
$\sigma$-algebra even though each finite-horizon restriction satisfies
$P_t \ll Q_t$.
\end{proposition}

\begin{proof}
Since $Z_t \geq 0$ is a $Q$-martingale, $Z_t \to Z_\infty$ $Q$-a.s.\ by
the almost-sure convergence theorem.  In case (ii), uniform integrability
gives $L^1(Q)$ convergence, hence $\E_Q[Z_\infty] = \lim_t \E_Q[Z_t] = 1$,
which is the condition for $P \ll Q$.  In case (iii), $\E_Q[Z_\infty] = 0$
forces $Z_\infty = 0$ $Q$-a.s.; one then constructs a set of full $P$-measure
and zero $Q$-measure using the set $\{Z_\infty = 0\}$.
\end{proof}

\begin{remark}[Sufficient condition via energy]
A sufficient condition for uniform integrability of $Z_t$ is
\begin{equation}
\label{eq:theta-square}
\int_0^\infty \theta_s^2\, ds < \infty \quad Q\text{-a.s.}
\end{equation}
Intuitively, \eqref{eq:theta-square} says the two measures differ only by a
finite total ``information'' over the infinite horizon.  When the drift
$\theta$ is stationary and non-zero, condition~\eqref{eq:theta-square}
fails, $P$ and $Q$ are mutually singular at infinity, and there is no
merging.
\end{remark}

\subsection{The Blackwell--Dubins theorem in continuous time}

\begin{theorem}[Merging on $[0, +\infty)$]
\label{thm:BD-CT}
Let $P$ and $Q$ be probability measures on $(\Omega, \mathscr{F}_\infty)$
with $P \ll Q$.  Let $Z_t = dP_t/dQ_t$ be the likelihood-ratio process.
Define
\[
d_t = \TV{P(\cdot\,|\,\mathscr{F}_t) - Q(\cdot\,|\,\mathscr{F}_t)}.
\]
Then $d_t \to 0$ almost surely under $Q$ (and hence under $P$) as
$t \to +\infty$.
\end{theorem}

\begin{proof}[Proof sketch]
The argument parallels the discrete-time proof.  Under $P \ll Q$, the
process $Z_t = \Ec{Z_\infty}{\mathscr{F}_t}_Q$ is a uniformly integrable
$Q$-martingale and $Z_t \to Z_\infty$ in $L^1(Q)$.  The continuous-time
analogue of \eqref{eq:supermg} holds by the same conditional-convexity
argument: $\{d_t, \mathscr{F}_t\}$ is a non-negative $Q$-supermartingale
bounded in $[0,1]$, hence $d_t \to d_\infty$ $Q$-a.s.  The $L^1$ bound
$\E_Q[d_t] \leq \tfrac{1}{2}\E_Q[|Z_t - Z_\infty|] \to 0$
forces $d_\infty = 0$ $Q$-a.s.
\end{proof}

\subsection{The role of the energy condition on path space}

On the Wiener space with drift $\theta$, the condition $P \ll Q$ is
equivalent to
\begin{equation}
\label{eq:Girsanov-condition}
\E_Q\!\left[\exp\!\left(\frac{1}{2}\int_0^\infty \theta_s^2\, ds\right)\right]
< \infty
\end{equation}
(Novikov on the infinite horizon), or equivalently to the stochastic
exponential $Z_t$ being uniformly integrable.

When \eqref{eq:Girsanov-condition} holds:
\begin{itemize}
\item Both $P$ and $Q$ assign probability one to paths on which the empirical
      average $\frac{1}{t}\int_0^t \theta_s\, ds \to 0$ as $t \to \infty$.
\item The observed increment $dW_t = dW_t^Q$ distinguishes the drift only
      locally, and the accumulated evidence is not sufficient to render the
      two measures singular.
\item By Theorem~\ref{thm:BD-CT}, the conditional distributions of the
      future path merge.
\end{itemize}
When \eqref{eq:Girsanov-condition} fails (e.g., $\theta_s = \theta_0 > 0$
constant), the sample-path average $\frac{1}{t}\int_0^t dW_s$ satisfies a
strong law under each measure that distinguishes them with probability one:
$P$ and $Q$ are mutually singular, and the conditional distributions
diverge rather than merge.

%----------------------------------------------------------------------
\section{Kakutani's Theorem and the Absolute-Continuity Condition}
\label{sec:kakutani}
%----------------------------------------------------------------------

A natural question is: when does the hypothesis of Theorem~\ref{thm:BD}
hold?  For product measures, the answer is given by a classical theorem of
\citet{kakutani1948}.

\subsection{Infinite product measures}

Let $P = \bigotimes_{n=1}^\infty P_n$ and $Q = \bigotimes_{n=1}^\infty Q_n$
be infinite product measures, where $P_n$ and $Q_n$ are probability measures
on $(S, \mathscr{S})$.  The \textbf{Hellinger affinity} between $P_n$ and
$Q_n$ is
\[
\rho_n = \int_S \sqrt{\frac{dP_n}{d\lambda}\cdot\frac{dQ_n}{d\lambda}}\, d\lambda,
\]
where $\lambda$ is any common dominating measure; $\rho_n \in (0, 1]$ with
$\rho_n = 1$ iff $P_n = Q_n$.

\begin{theorem}[Kakutani's Dichotomy, 1948]
\label{thm:kakutani}
For infinite product measures:
\[
P \sim Q \quad \iff \quad \prod_{n=1}^\infty \rho_n > 0
\quad \iff \quad \sum_{n=1}^\infty (1 - \rho_n) < \infty.
\]
If $\prod_{n=1}^\infty \rho_n = 0$, then $P \perp Q$.  There is no
intermediate case: either the measures are equivalent or mutually singular.
\end{theorem}

\begin{proof}[Proof sketch]
The likelihood ratio $Z_n = \prod_{k=1}^n (dP_k/dQ_k)$ is a $Q$-martingale.
The product $\E_Q[\sqrt{Z_n}] = \prod_{k=1}^n \rho_k$ is the Hellinger
integral.  If $\prod \rho_k > 0$, then $Z_n$ is bounded in $L^{1/2}$, which
implies uniform integrability in $L^1$ and hence $P \ll Q$.  If
$\prod \rho_k = 0$, then $Z_n \to 0$ in $L^{1/2}$, which forces
$Z_\infty = 0$ $Q$-a.s., establishing $P \perp Q$.
\end{proof}

\subsection{Implication for merging}

Kakutani's theorem identifies precisely the regime in which Blackwell--Dubins
applies to i.i.d.-type sequences.  If the one-period measures $P_n$ and $Q_n$
are ``close'' in the sense that $\sum_n (1 - \rho_n) < \infty$, then merging
is guaranteed.  If they differ by a fixed total variation amount at every
period (e.g., $P_n$ and $Q_n$ are distinct fixed measures independent of $n$),
the Hellinger product converges to zero and the two priors are mutually
singular: Blackwell--Dubins does not apply, and the agents never agree.

%----------------------------------------------------------------------
\section{Applications}
\label{sec:applications}
%----------------------------------------------------------------------

\subsection{Bayesian learning}

The most direct application is to Bayesian inference.  Suppose the data
$(x_1, x_2, \ldots)$ are drawn from the true measure $Q^*$.  An agent holds
a prior $\pi$ over a family $\{Q_\theta : \theta \in \Theta\}$ of models,
inducing a marginal $P = \int Q_\theta\, \pi(d\theta)$.  If $P \ll Q^*$ (the
true model has positive probability under the prior), then by
Theorem~\ref{thm:BD}, the agent's posterior predictive distribution merges
with $Q^*$ in total variation:
\[
\TV{P(\cdot\,|\,x_1, \ldots, x_n) - Q^*(\cdot\,|\,x_1, \ldots, x_n)} \to 0
\quad Q^*\text{-a.s.}
\]
This is a strong form of Bayesian consistency: the agent's predictions about
all future events become arbitrarily accurate, regardless of the specific
prior, as long as the prior assigns positive probability to the truth.

\subsection{Rational expectations and the common prior assumption}

In macroeconomics, the merging-of-opinions theorem provides a foundation for
the common prior assumption embedded in rational expectations models.  If two
agents $P$ and $Q$ have mutually absolutely continuous priors and observe a
common history, Theorem~\ref{thm:BD} guarantees that they will eventually
agree on all conditional forecasts.

Conversely, \citet{aumann1976}'s agreement theorem shows that agents with
a common prior who share their posteriors cannot ``agree to disagree'' on
conditional probabilities.  Blackwell--Dubins complements this by showing
that even with heterogeneous priors, merging occurs eventually if the priors
are in the same absolute-continuity class.

\subsection{Signal extraction and diffusion models}

Consider the continuous-time signal-extraction model
\[
dy_t = \mu_t\, dt + \sigma\, dW_t,
\]
where $\mu_t$ is an unobserved signal process and $W_t$ is noise.  Two
agents with different priors over $\{\mu_t\}$ observe the same path
$\{y_s : s \leq t\}$.  Their respective filtering equations
(Kalman--Bucy or Wonham) produce conditional distributions $P(\mu_t|y^t)$
and $Q(\mu_t|y^t)$.

By Theorem~\ref{thm:BD-CT}, if their priors over the signal process are
mutually absolutely continuous, their filtered distributions merge as
$t \to \infty$.  The rate of merging is governed by the signal-to-noise
ratio: when $\sigma$ is small, the signal is identified quickly and merging
is fast; when $\sigma$ is large, merging is slow.

\subsection{Ergodic Markov chains}

For a Markov chain with transition kernel $\Pi$ and two initial distributions
$P_0 = \mu$ and $Q_0 = \nu$, the $n$-step distributions are $\mu \Pi^n$ and
$\nu \Pi^n$.  If $\Pi$ is ergodic with unique stationary distribution $\pi$,
then both $\mu \Pi^n$ and $\nu \Pi^n$ converge to $\pi$, and in particular
\[
\TV{\mu\Pi^n - \nu\Pi^n} \leq \TV{\mu\Pi^n - \pi} + \TV{\nu\Pi^n - \pi} \to 0.
\]
This is a version of the merging result that does \emph{not} require absolute
continuity, because the ergodicity of $\Pi$ already forces convergence to a
common limit.  The Blackwell--Dubins theorem is the appropriate generalisation
for non-Markovian or non-ergodic processes, where no single invariant measure
exists but absolute continuity of the priors is the operative condition.

%----------------------------------------------------------------------
\section{Related Results and Extensions}
\label{sec:related}
%----------------------------------------------------------------------

\subsection{Lévy's zero-one law and Doob's consistency theorem}

\citet{levy1937} established that for any event $A \in \mathscr{F}$,
\[
\Ec{\mathbf{1}_A}{\mathscr{F}_n} \to \mathbf{1}_A \quad \text{a.s.\ and in } L^1.
\]
This is a precursor to Blackwell--Dubins: it says that \emph{under a single
measure}, the conditional probability of any fixed event concentrates on 0
or 1 as information grows.  Blackwell--Dubins strengthens this to compare
conditional probabilities across two \emph{different} measures.

\citet{doob1949} proved that, under regularity conditions, the posterior
distribution of a parameter $\theta$ given data $x_1, \ldots, x_n$ converges
to the point mass at the true $\theta^*$ for almost every $\theta^*$ under
the prior.  This ``Doob consistency'' requires the prior to assign positive
probability to neighbourhoods of $\theta^*$, which is a form of absolute
continuity.

\subsection{Diaconis--Freedman and robustness of Bayesian consistency}

\citet{diaconis1986} study conditions under which Bayesian procedures are
consistent in infinite-dimensional settings.  They show that absolute
continuity of the prior with respect to the truth is necessary for Doob
consistency, connecting their analysis directly to the Blackwell--Dubins
framework.  They also give examples of priors that are inconsistent when the
truth lies outside the support of the prior---a setting in which $P \perp Q$
and Blackwell--Dubins does not apply.

\subsection{Rate of merging}

Blackwell--Dubins establishes only almost-sure convergence, not a rate.
Rates of merging depend on the ``distance'' between $P$ and $Q$ in a suitable
sense.  For parametric Bayesian models, the posterior contracts at the
$1/\sqrt{n}$ rate (Bernstein-von Mises theorem).  For nonparametric models,
rates are slower.  A useful bound relates the rate of $L^1$ convergence of
$Z_n$ to the rate of merging:
\[
\E_Q[d_n] \leq \tfrac{1}{2}\E_Q[|Z_n - Z_\infty|],
\]
so any rate for the right-hand side transfers to $d_n$.

\subsection{Merging without absolute continuity}

\citet{lehrer1996} and subsequent work in game theory have explored
conditions weaker than absolute continuity that still imply merging of
predictions on some subsequence of times.  These extensions are relevant for
the analysis of heterogeneous-agent models in which agents' priors need not
be mutually absolutely continuous.

\subsection{Summary of the main logical dependencies}

The chain of implications underlying the Blackwell--Dubins theorem can be
summarised as follows:
\begin{align*}
P \ll Q
\;&\Longrightarrow\;
Z_n = \mathbb{E}_Q[Z_\infty\,|\,\mathscr{F}_n] \text{ uniformly integrable under } Q \\
\;&\Longrightarrow\;
Z_n \xrightarrow{L^1(Q)} Z_\infty
\;\Longrightarrow\;
d_n \xrightarrow{L^1(Q)} 0
\;\Longrightarrow\;
d_n \xrightarrow{Q\text{-a.s.}} 0.
\end{align*}
The first implication is the Radon--Nikodym structure.  The second is
the $L^1$ martingale convergence theorem.  The third is the TV-distance
bound~\eqref{eq:dn-L1} together with the supermartingale property.  The
fourth combines the a.s.\ convergence of the supermartingale $d_n$ with the
$L^1$ convergence to pin the a.s.\ limit at zero.

\bigskip
\bibliographystyle{plainnat}
\bibliography{blackwell_dubins}

\end{document}
